\SourceLeaf{379}%
\needspace{12\baselineskip}
\begin{center}[3.]\end{center}

The most general solution to the problem of determining all the shapes that a given curved surface can assume rests on the following. Let $P$ be the point on the celestial sphere corresponding to the normal to the surface of the superficies in one shape, and $Q$ the point corresponding to the same point of the surface in another shape. Likewise, the points $P',P'',P'''$ etc. for one shape correspond to the points $Q',Q'',Q'''$ [etc.] for the other. The solution then consists in requiring that every closed figure $PP'P''P'''\, .\, .\, .$ on the sphere have the same area as the corresponding figure $QQ'Q''Q'''\, .\, .\, .$.

\textls[120]{Later note.} This beautiful theorem is entirely correct, but it is insufficient for the solution.

\needspace{8\baselineskip}
\begin{center}REMARKS.\end{center}

Notes [2] and [3] are found in the notebook “Den astronomischen Wissenschaften gewidmet,” which \textsc{Gauss} began in November 1801. They presumably date from about the same time as the letter from \textsc{Gauss} to \textsc{Schumacher} printed under [1].
\begin{flushright}\textsc{Stäckel}.\end{flushright}
